\reviewsection{The Environment Determines Which Competence Becomes Visible}

\citet{downing2015classroom} extend this lesson into the general education classroom by asking educators to examine actual routines, exchanges, peer relationships, and communication opportunities. A student can be physically present and still have little access to authorship. Predictable routines may help a learner anticipate an exchange, but meaningful participation also requires vocabulary and partner support for unpredictable comments, questions, disagreement, humor, and repair. Communication cannot be reduced to selecting between two adult-approved choices.

This reading connects directly to my K--5 technology teaching placement. I saw participation change when I reduced language demands, modeled the exact technology steps, made the sequence visible, and allowed familiar tools to remain available. The learner's competence had not suddenly appeared; the conditions for showing it had changed. That experience reminds me that implementation belongs to the environment and communication partners, not only to the student.

My EDU486 Invisible Invaders camp work gives me another concrete picture of this principle. Youth could contribute by looking, moving, choosing, handling tools, photographing, drawing, writing, explaining, partnering, pausing, or returning. Identity Beads represented different and complementary contributions rather than ranking children by one performance. During evidence-to-action work, youth revised models with labels and arrows and could create wearable advocacy messages with words, images, or symbols. At the showcase, teaching could happen through demonstrating, pointing, answering, asking, holding evidence, partnering, or observing \citep{campEvidence2026}.

Those routes were not all AAC, and I am not claiming that a particular camper used an AAC system. They do show how classroom design decides which forms of authorship become visible. A communication-rich classroom distributes ways to enter the work before a student is treated as disengaged, incapable, or unwilling.
